\documentclass[leqno, a4paper, 12pt]{jsarticle}
\usepackage{amssymb}
\usepackage{amsthm}
\usepackage{amsmath}
\usepackage{comment}
\usepackage{url}

\usepackage{enumitem}
%\usepackage{showkeys}


%定理環境 thmはあまり使わずpropをなるべく使う。
%主張は基本的に証明の中で使い、そのため番号を付けない。
%注意環境を付けてみた。
\newtheorem{thm}{定理}
\newtheorem{prop}[thm]{命題}
\newtheorem{cor}[thm]{系}
\newtheorem{lem}[thm]{補題}
\newtheorem{df}[thm]{定義}
\newtheorem{exam}[thm]{例}
\newtheorem{fact}[thm]{事実}
\newtheorem*{claim}{主張}
\newtheorem*{cau}{注意}
\newtheorem*{rmk}{注意}
\renewcommand{\proofname}{証明}
%矢印たち。
%\toは\rightarrowと同じ。\getsは\leftarrowと同じ。同値は\iff
\newcommand{\To}{\Longrightarrow}
\newcommand{\Gets}{\LongLeftarrow}
%黒板太字
\newcommand{\nn}{\mathbb{N}}
\newcommand{\zz}{\mathbb{Z}}
\newcommand{\qq}{\mathbb{Q}}
\newcommand{\rr}{\mathbb{R}}
\newcommand{\cc}{\mathbb{C}}
%無限大
\newcommand{\mgn}{\infty}
%差集合は\setminusで統一する。等号の否定は\neq
%真部分集合は\subsetneqq　偏微分は\partial
%ディスプレイスタイル。\dfracでディスプレイ分数
\newcommand{\dpst}{\displaystyle}

\title{論文メモ
\large{\\--- Noble? ---}
}
\author{ナンブ　キトラ}
\date{\today}
\begin{document}
\maketitle


本棚を整理したいので，
溜まっている論文のメモを書いて未来への手紙とすることにする．


以下の論文は，
位相空間
$X$について
直積
$X^{\mathfrak{m}}$
が任意の基数
$\mathfrak{m}$
について
正規ハウスドルフならば
$X$
はコンパクトハウスドルフ
というNoble
の定理から発展した論文を集めた論文の束だと思う．
このNobleの定理については
電波通信でも紹介している
(https://concious4410.hatenablog.com/entry/2016/07/06/223459)．
証明の方針としては，
$X$が一点ならば話は簡単なので$X$は2点以上含んでいるとしても良い．
まずStoneの定理というのがあり，
これはコンパクトでない距離空間$Z$と
非可算濃度
$\mathfrak{m}
$については
$Z^{\mathfrak{m}}$
は正規ではないという定理である．
ちなみにこの定理を使うと完全正則だが正規でない例がたくさん作れる．
このことから，もしも位相空間
$X$について$X^{\mathfrak{m}}$
が任意の基数
$\mathfrak{m}$
について正規ならば
$X$が可算コンパクトであることがわかる．
これは背理法によって示される．
$X$に閉集合であるような
可算離散空間
$D$が存在すると
$D^{\mathfrak{m}}$
は
正規空間
$X$の閉部分空間なので正規でなければならないが，
これはStoneの定理に矛盾というわけである．
また森田の
定理というのがありそれは
$X\times \{0, 1\}^{\mathfrak{m}}$
が正規ならば
$X$
は$\mathfrak{m}$-パラコンパクトというものである．
空間$X$が2点以上で$X^{\mathfrak{m}}$
がいつも正規ならばこの条件を満たすので
$X$は任意の$\mathfrak{m}$について
$\mathfrak{m}$-パラコンパクトである．つまりパラコンパクトである．
さてパラコンパクトな可算コンパクト空間はコンパクトなので
$X$
はコンパクトハウスドルフである．

まあそれはさておき
Nobleの定理の簡単な証明に関するものは
\cite{engelking1988elementary}, 
\cite{franklin1972normality}, 
\cite{keesling1972normality}, 
\cite{polkowski1979n}
である．


論文
\cite{noble1967generalization}
はNoble自身による論文でStoneの定理に関するものである．
これは上の
Nobleの定理の萌芽だと思う．

Nobleの原論文は
\cite{MR0250261}
と
\cite{MR0283749}
である．
論文?
\cite{Fleischer1967}
も参照のこと．


Noble
の定理に関連する概念として
$[\mathfrak{a}, \mathfrak{b}]$-compactness
がある．これは別にNobleの定理以前からあるやつだと思う．
またこの概念は
M. E. RudinによるDowker spaceの構成にも関係ある．
これに関連する論文は以下の通りである．
\cite{Vaughan1975}, 
\cite{LIPPARINI20061365}, 
\cite{butcherjoseph1976}, 
\cite{HODEL1974179}, 
\cite{shirayama2007note}, 
\cite{gal1958theory}. 



%READ!
%myplain is the same to amsplain style. 
\nocite{*}
\bibliographystyle{myplaindoidoi}
\bibliography{noble.bib}



\end{document}